\chapter{Strategy Taxonomy}
\label{app:taxonomy}

The final six-category taxonomy (v4) constrains the \techil{<plan>} block
to the following therapeutic rescripting mechanisms. Each category describes
\emph{how} the nightmare changes, not what outcome it achieves. The
taxonomy evolved through four iterations to resolve distribution skew
(\autoref{sec:meth-taxonomy}).

\begin{table}[h]
  \centering
  \caption{Strategy taxonomy v4: full definitions as provided to the
    therapist model. NOT-examples reduce category overlap.}
  \label{tab:taxonomy-full}
  \small
  \begin{tabularx}{\textwidth}{lX}
    \toprule
    Category & Definition (as given to the model) \\
    \midrule
    \techil{confrontation} &
      The dreamer directly faces, fights, or stands up to the threat itself
      (NOT avoiding, escaping, or gaining general confidence). \\
    \techil{self\_empowerment} &
      The dreamer gains a new ability, skill, or inner strength they did
      not have before (\eg discovers they can fly, becomes stronger, unlocks
      a hidden power). The DREAMER changes, not the threat. (NOT
      reinterpreting existing elements, which is \techil{cognitive\_reframe}.) \\
    \techil{safety} &
      A physical barrier, shelter, or protective figure is added to shield
      the dreamer (\eg walls, locked doors, a guardian). NOT just feeling
      calm or changing the scene. \\
    \techil{cognitive\_reframe} &
      The meaning or identity of existing dream elements changes (\eg a
      monster becomes harmless, an enemy becomes a friend). The THREAT
      changes, not the dreamer. (NOT the dreamer gaining new powers, which
      is \techil{self\_empowerment}.) \\
    \techil{social\_support} &
      People in the dream become supportive, or a new ally is introduced.
      Existing characters can shift from threatening to encouraging. Any
      form of interpersonal support counts. \\
    \techil{sensory\_modulation} &
      The dream's sensory qualities change: what is seen, heard, or felt
      shifts (\eg darkness becomes light, harsh sounds become soft music,
      the dreamer feels calm through breathing or self-talk). Covers both
      external environment changes and internal somatic calming. \\
    \bottomrule
  \end{tabularx}
\end{table}

\paragraph{Revision history.}
\begin{description}
  \item[v1 (8 categories):] Included \techil{empowerment} and
    \techil{mastery} as separate categories. These two consumed 87.7\% of
    all strategy picks across 1{,}272 trials, while three categories were
    never selected. Root cause: both describe a therapeutic \emph{goal}
    (the patient gains control) rather than a \emph{mechanism}.
  \item[v2 (7 categories):] Merged \techil{empowerment} and
    \techil{mastery} into \techil{agency}. The merged category achieved
    100\% dominance.
  \item[v3 (7 categories):] Split \techil{agency} along the mechanism
    dimension into \techil{confrontation} (external action) and
    \techil{self\_empowerment} (internal transformation). Sharpened all
    descriptions to focus on mechanism, not outcome.
  \item[v4 (6 categories):] Merged \techil{emotional\_regulation} into
    \techil{sensory\_modulation}. In smoke testing (120 trials),
    \techil{emotional\_regulation} appeared once (0.4\%). In single-turn
    rescripting, internal calming manifests as sensory change.
\end{description}


\chapter{Prompts}
\label{app:prompts}

This appendix reproduces the key prompts used in the evaluation pipeline.
All prompts are stored as YAML files in \techil{data/prompts/}.

\section{Fused Plan and Response Prompt}
\label{app:fused-prompt}

The following system prompt is injected into the therapist model during
evaluation. The \techil{\{categories\_block\}} placeholder is replaced
with the full taxonomy definitions from \autoref{tab:taxonomy-full}.

\begin{small}
\begin{verbatim}
You are an Imagery Rehearsal Therapy (IRT) assistant
helping a patient rescript their nightmare.

First, select 1-2 therapeutic strategies from this taxonomy:
{categories_block}

Choose the strategies most specifically suited to THIS
patient's nightmare content. Avoid defaulting to general
approaches when a more specific mechanism fits.

Then respond to the patient using those strategies.

STRICT output format:
Line 1: <plan>category1 / category2</plan>
Line 2: (empty)
Line 3+: Your therapeutic response (conversational,
         supportive, max 3 sentences)
\end{verbatim}
\end{small}

\section{Alignment Judge Prompt}
\label{app:judge-prompt}

The following prompt is sent to \modelname{Gemini 3.1 Pro} ($T = 1.0$,
thinking enabled). The \techil{\{taxonomy\_block\}} is replaced with
the full strategy definitions. The \techil{\{strategies\_block\}} and
\techil{\{response\}} are filled per trial.

\begin{small}
\begin{verbatim}
System:
You are evaluating whether a therapist's response implements
specific therapeutic strategies from IRT.

## Strategy Definitions
{taxonomy_block}

## Scoring Rubric
For each declared strategy, assign ONE score:
- 2 (implemented): The response clearly and specifically
  implements this strategy.
- 1 (partial): The response touches on this strategy but
  does not fully develop it.
- 0 (absent): The response does not implement this strategy.

## Instructions
1. Read the therapist response carefully.
2. For each declared strategy, provide brief reasoning
   (1-2 sentences), then your score.
3. Use EXACTLY this output format for each strategy:

strategy_id: <reasoning> | score: <0|1|2>

User:
## Declared Strategies
{strategies_block}

## Therapist Response
{response}
\end{verbatim}
\end{small}

\section{Example Patient Vignette}
\label{app:vignette-example}

Six vignettes define patient profiles for dialogue generation. The
following example shows the \emph{anxious} profile. All vignettes follow
the same schema.

\begin{small}
\begin{verbatim}
{
  "type": "anxious",
  "name": "Maya",
  "age": 24,
  "gender": "female",
  "background": "A graduate student dealing with nightmares
    that started during a particularly stressful thesis
    period. Tends to catastrophize and worry about whether
    things will ever get better.",
  "nightmare": {
    "content": "Being in an exam room where the questions
      keep changing, the clock is speeding up, and everyone
      else has finished. Sometimes the room starts flooding
      or the walls close in.",
    "frequency": "5-6 times per week",
    "duration": "8 months"
  },
  "personality_traits": [
    "worried", "seeking reassurance",
    "catastrophizing", "hyperaware of sensations",
    "needs frequent validation"
  ],
  "resistance_level": "moderate",
  "sample_responses": [
    "What if I try to change the dream and it makes it
     worse? Has that ever happened?",
    "I'm already anxious just thinking about the nightmare.
     I don't know if I can do this."
  ]
}
\end{verbatim}
\end{small}


\chapter{Supplementary Results}
\label{app:results}

This appendix contains figures that support the main text but were deferred
to keep the results chapter focused on the primary findings.

\section{Plan Validity Rate}

\begin{figure}[h]
  \centering
  \includegraphics[width=0.85\textwidth]{figures/fig_A1_validity.pdf}
  \caption{Plan validity rate by model, pooled across temperatures and
    vignettes. All models exceed 99\%. The y-axis begins at 0.95 to show
    the small differences between models.}
  \label{fig:app-validity}
\end{figure}

\section{BERTScore by Model and Temperature}

\begin{figure}[h]
  \centering
  \includegraphics[width=\textwidth]{figures/fig_5_3_bertscore.pdf}
  \caption{Mean BERTScore F1 by model and temperature, grouped by family.
    The pooled temperature effect is stronger than for Jaccard
    ($\rho = -0.486$), but BERTScore is insensitive to plan-level
    differences (\autoref{sec:res-similarity}).}
  \label{fig:app-bertscore}
\end{figure}

\section{Vignette Heatmaps}

\begin{figure}[h]
  \centering
  \includegraphics[width=\textwidth]{figures/fig_5_4_vignette_slice.pdf}
  \caption{Jaccard by model and vignette, one panel per temperature.
    Some vignettes (\eg resistant, trauma) show lower consistency across
    multiple models. The significant vignette effect ($p = 0.007$) is
    visible as column-wise variation.}
  \label{fig:app-vignette}
\end{figure}

\section{Slice Depth Heatmap}

\begin{figure}[h]
  \centering
  \includegraphics[width=0.85\textwidth]{figures/fig_slice_depth_heatmap.pdf}
  \caption{Jaccard by vignette and conversation depth for
    \modelname{Mistral Large~3}, one panel per temperature. Slice~2 shows
    the lowest consistency across most vignettes, supporting its selection
    as a conservative evaluation point.}
  \label{fig:app-slice-heatmap}
\end{figure}
